% =====================================================================
%  c26 —— 原文印刷页 211--222（PDF 第 212--223 页）
%  内容：附录 I（I. Chern-Weil Theory And Chern-Simons Forms）
%        I.1 特征类；I.2 基本形式；I.3 Lie 代数上的不变多项式（I.3.1 例子）；
%        I.4 Chern-Weil 同态；I.5 对联络的依赖：特征类的构造；
%        I.6 Borel 定理；I.7 次级特征类：Chern-Simons 形式
%  编号公式：I.1--I.54 共 54 条（多行者为 I.5、I.16、I.19、I.20、I.22、
%        I.33、I.42、I.46、I.47、I.49、I.52）。
%  插图：Figure 37（第 220 页）、Figure 38（第 221 页）。原 PDF 中两图均为
%        空白占位框（figures/fig37.png、fig38.png 亦为空白），故按 main.tex
%        中既定的 \MissingFigureNote 方案处理：保留图号与图注，不插入空白图。
%  原文印刷页 211 上半页为附录 H 习题（H.36--H.40）的续页，属上一片段范围，
%  本片段按约定自 \section 级附录 I 开始。
% =====================================================================

\section{Chern-Weil 理论与 Chern-Simons 形式} % 原文: I. Chern-Weil Theory And Chern-Simons Forms

\begin{keytakeaways}
  \item 特征类是主丛同构类的函数：$\xi$ 把以 $G$ 为结构群的主丛 $P\to M$ 送到整上同调类
    $\xi(P)\in H^*(M;\mathbb{Z})$，并要求自然的拉回性质 \eqref{eq:I.1}；等价地，它来自分类空间
    的上同调类 $\theta\in H^*(BG;\mathbb{Z})$ 沿分类映射的拉回 $\xi(P)=f^*(\theta)$。
  \item Chern-Weil 构造：给定联络 $\Theta$ 与 Lie 代数上的不变多项式 $\varphi$（次数 $k$），把曲率
    $\bar F\in\Omega^2(P;\mathfrak{g})$ 代入得 $2k$ 次形式 $\varphi(\bar F)=\pi^*(\bar\varphi)$；它是基本
    形式且闭，故给出分次环同态 $I^*(G)\to H^*_{DR}(M)$，即 Chern-Weil 同态。
  \item 与联络无关：沿联络空间中任一路径 $A(t)$ 积分，两个联络处的特征形式之差是恰当形式
    \eqref{eq:I.20}、\eqref{eq:I.23}，因此上同调类与联络选取无关；在联络族上引入双分次与万有
    联络后，这一机制被整理成下降方程 \eqref{eq:I.42}。
  \item 次级特征类：把上述恰当形式中的原函数取出，便得相对 Chern-Simons 形式
    $CS_{2n-1}(A_1,A_0)$，它满足 $dCS_{2n-1}=\operatorname{Tr}F_1^n-\operatorname{Tr}F_0^n$。它本质上是
    「两个联络之间」的量，只有在平凡丛上取 $A_0=0$ 时才写作单个联络的 $CS(A)$；其积分只
    定义到模 $\mathbb{Z}$。
  \item Borel 定理：$H^*(BG;\Lambda)=(H^*(BT,\Lambda))^W$ 与 $I^*(G)=S^*(\mathfrak{g}^*)^W$ 给出
    $I^*(G)\cong H^*(BG;\mathbb{R})$，把不变多项式环、Weyl 群不变量与分类空间上同调三者闭合起来，
    这是 Chern-Weil 同态「取遍所有特征类」的根本原因。
\end{keytakeaways}

\subsection{特征类} % 原文: I.1 Characteristic classes

我们没有足够的篇幅来充分展开这个重要的主题。下面是电报式的简述。

粗略地说，特征类（characteristic class）是度量向量丛与主丛的扭曲（twisting）程度的一种办法。

按定义，特征类是一条把上同调类指派给主丛（的同构类）的规则。当 $\pi: P\to M$ 是以 $G$ 为结构群的主丛时，特征类是 $M$ 上的整上同调类。因此特征类是从 $M$ 上主丛的同构类集合到 $H^*(M;\mathbb{Z})$ 的映射。若 $P\to M$ 是主丛，我们用 $\xi(P)$ 记相应的上同调类。此外，这个映射应当是\emph{自然的}，即
\begin{equation}
  f^*(\xi(P)) = \xi(f^*(P))
  \tag{I.1}
  \label{eq:I.1}
\end{equation}

注意这样的特征类构成一个环。

由于性质 \eqref{eq:I.1}，特征类可以用分类空间 $BG$ 的上同调类来定义。若 $\theta\in H^*(BG;\mathbb{Z})$，而 $f: M\to BG$ 是 $P$ 的分类映射，即 $P=f^*(EG)$，则可取 $\xi(P)=f^*(\theta)$。而且所有特征类都是以这种方式得到的。

当丛 $P$ 上配有一个联络时，我们可以构造特征类的 DeRham 代表元。尽管这样做会丢失挠（torsion）信息，这一构造仍然非常有用，在物理中也极为常见。它导致 Yang-Mills 作用量中的拓扑项，并且通过所谓的次级特征类（即 Chern-Simons 形式），引出几何与物理相互作用中最有趣、最重要的许多例子。

\subsection{基本形式} % 原文: I.2 Basic forms

\textbf{定义} 丛 $P\to M$ 上的一个\emph{基本形式}（basic form）是指由底空间上的形式 $\tilde{\omega}$ 拉回而得的形式：$\omega=\pi^*(\tilde{\omega})$。

\textbf{命题.} 设 $P$ 是紧致连通群 $G$ 的主 $G$ 丛。则 $\omega\in\Omega^*(P)$ 是基本形式，当且仅当它满足
\begin{enumerate}[label=(\alph*)]
  \item $\iota(\xi(X))(\omega)=0$；
  \item $\mathcal{L}(\xi(X))(\omega)=0$，
\end{enumerate}
其中 $X$ 取遍 $\mathfrak{g}$。

\textbf{证明}：局部地取坐标 $(x,g)$。于是总空间上的任何形式都可以用 $dx^\mu$ 与 $e^a$ 展开，其中 $\Theta=e^aT_a$。

垂直性方程 $\iota(\xi(X))(\omega)=0$ 表示 $\omega$ 没有垂直分量。也就是说，在局部坐标下它具有如下形式：
\begin{equation}
  \omega = \omega_{\mu_1\cdots\mu_k}(x,g)\,dx^{\mu_1\cdots\mu_k}
  \tag{I.2}
  \label{eq:I.2}
\end{equation}

其次，Lie 导数条件表明系数函数必须与 $g$ 无关。\ $\clubsuit$

% [Tutor Note: 这里被略去的一步是：由 b) 得 $\mathcal{L}(\xi(X))\omega=0$，而 $\mathcal{L}=\iota d+d\iota$，把 (I.2) 代入并对 $g$ 沿单参数子群求导，即得 $\partial_g\omega_{\mu_1\cdots\mu_k}=0$（对连通 $G$ 而言 $g$ 可沿生成元方向走到单位元）。这正是不必额外假设「与 $g$ 无关」的原因，也说明了命题为何要假设 $G$ 紧致\emph{且连通}：对不连通群，(b) 只能保证在单位元分支上不变，故须换成 (I.3) 的 $R_g^*(\omega)=\omega$。]

\textbf{注}：更一般地，若 $G$ 不连通，则上面的条件 (b) 应表述为如下条件：
\begin{equation}
  R_g^*(\omega)=\omega \qquad \forall g\in G
  \tag{I.3}
  \label{eq:I.3}
\end{equation}

\subsection{Lie 代数上的不变多项式} % 原文: I.3 Invariant polynomials on the Lie algebra

设 $\mathfrak{g}$ 是 $\mathbb{F}=\mathbb{R}$ 或 $\mathbb{C}$ 上的 Lie 代数。$\mathfrak{g}$ 上的一个 $k$ 次多项式指的是一个全对称的多重线性函数：
\begin{equation}
  \varphi:\mathfrak{g}^{\otimes k}\to\mathbb{F}
  \tag{I.4}
  \label{eq:I.4}
\end{equation}

对称多项式是 $\mathrm{Sym}^k\mathfrak{g}^*$ 的元素，而 $\oplus_k\mathrm{Sym}^k\mathfrak{g}^*$ 合起来构成一个代数。次数分别为 $k_1,k_2$ 的两个多项式 $\varphi_1$ 与 $\varphi_2$ 的乘积是
\begin{equation}
\begin{aligned}
  (\varphi_1\cdot\varphi_2)(v_1,\dots,v_{k_1+k_2}) :={}&
    \frac{1}{(k_1+k_2)!}\sum_{\sigma\in S_{k_1+k_2}}
    \varphi_1(v_{\sigma(1)},\dots,v_{\sigma(k_1)})\\
  &\times \varphi_2(v_{\sigma(k_1+1)},\dots,v_{\sigma(k_1+k_2)})
\end{aligned}
  \tag{I.5}
  \label{eq:I.5}
\end{equation}

% [Tutor Note: 「$k$ 次多项式」在本文中始终指\emph{齐次}的对称多重线性函数，即 $\mathrm{Sym}^k\mathfrak{g}^*$ 中的元素，而不是通常意义下总次数为 $k$ 的多项式。核对 (I.5) 即知：两个齐次多项式之积仍是齐次的（次数相加），系数 $1/(k_1+k_2)!$ 正是把「有序的 $(k_1+k_2)$ 元组」对称化所需的归一化。作者在 (I.9) 处又用 "homogeneous of degree $k$" 这一说法，两者一致；若不注意这一点，(I.5) 前面的因子极易漏写。]

一个\emph{不变多项式}（invariant polynomial）是指 $Ad(G)$ 不变的多项式。不变多项式构成 $S^*(\mathfrak{g}^*)$ 的一个子环，记作 $I(G)^*$ 或 $\mathrm{Inv}_G S^*(\mathfrak{g})$。

\textbf{注}：不变多项式完全由它在对角线上的取值 $\varphi(v,v,\dots,v)$ 决定。证明：取一组基，用未定元 $x_i$ 把 $v=\sum v_ix_i$ 代入。我们将滥用记号，把这个对角值也记作 $\varphi(v)$。于是 $\varphi(gvg^{-1})=\varphi(v)$。

\textbf{定理} 若 $G$ 是紧致单 Lie 群，则不变多项式环是 Casimir 元上的多项式环
\begin{equation}
  \mathrm{Inv}_G S^*(\mathfrak{g}) \cong \mathbb{R}[C_1,C_2,\dots,C_\ell]
  \tag{I.6}
  \label{eq:I.6}
\end{equation}
其中 $\ell$ 是 $G$ 的秩，而 $C_i$ 的次数为 $k_i$，即 $G$ 的指数（exponents）。

\textbf{证明}：不变多项式由它在对角线上的值 $\varphi(v)$ 决定，而由 $Ad$-不变性，它又由在一个 Cartan 子代数 $\mathfrak{t}\subset\mathfrak{g}$ 上的取值决定。反之，$\mathfrak{t}$ 上任何 $W$-不变的多项式都定义了 $\mathfrak{g}$ 上的一个不变多项式，这里 $W$ 是 Weyl 群。由于 $W$ 是有限群，我们总可以把 $\mathfrak{t}$ 上的多项式取平均，从而得到一个 $W$-不变的多项式。结果表明，$W$-不变多项式环本身也是一个多项式环：见 C. Chevalley, ``Invariants of finite groups generated by reflections,'' Amer. J. Math. \textbf{77}(1955)778。\ $\clubsuit$

% [Tutor Note: 证明中 "It turns out that ..." 一句省略的是 Chevalley 定理（有限反射群的不变多项式环是多项式环，生成元的次数等于该群的指数）；这里 $W$ 是 $G$ 的 Weyl 群，它在 $\mathfrak{t}$ 上由单反射生成，故生成元 $C_i$ 的次数 $k_i$ 正是 $G$ 的指数。若要补全，可查 Humphreys《Reflection Groups and Coxeter Groups》第 3 章，或 Bourbaki《Lie Groups and Lie Algebras》第 V 章 §5。作者此处只给了一篇原始文献。]

\subsubsection{例子} % 原文: I.3.1 Examples

例 1：$Lie(GL(n,\mathbb{C}))=Mat_n(\mathbb{C})$。不变多项式环由 Chern 多项式 $C_k$ 生成，后者由下式定义：
\begin{equation}
  \det\left(xI+\frac{1}{2\pi i}A\right) = \sum_j C_j(A,\dots,A)\,x^{n-j}
  \tag{I.7}
  \label{eq:I.7}
\end{equation}

例 2：GL(n,R)

例 3：U(N)

例 4：$G=SO(2r)$……

对 $G=SO(2r)$，会出现一个新的重要不变多项式——Pfaffian：
\begin{equation}
  \mathrm{Pfaff}(m) = \frac{1}{2^r r!}\sum_{\sigma\in S_{2r}}\mathrm{sign}(\sigma)\,
  m_{\sigma(1)\sigma(2)}m_{\sigma(3)\sigma(4)}\cdots m_{\sigma(2r-1)\sigma(2r)}
  \tag{I.8}
  \label{eq:I.8}
\end{equation}
对于反对称矩阵而言，它是矩阵元的一个多项式，且满足 $(\mathrm{Pfaff}(m))^2=\det m$。

例 5：$G=SO(2r+1)$。

例 6：$G=USp(2r)$。

\paragraph{练习}
证明对任意 $g\in O(2r)$ 有 $\mathrm{Pfaff}(gmg^{-1})=\det(g)\,\mathrm{Pfaff}(m)$。

\subsection{Chern-Weil 同态} % 原文: I.4 The Chern-Weil homomorphism

现在设 $\pi: P\to M$ 是以 $G$ 为结构群的主丛，配有联络 $\Theta$ 与场强 $\bar F\in\Omega^2(P;\mathfrak{g})$。设 $\varphi$ 是 $\mathfrak{g}$ 上的一个不变多项式。

于是我们可以考虑
\begin{equation}
  \varphi(\bar F)
  \tag{I.9}
  \label{eq:I.9}
\end{equation}
若 $\varphi$ 是 $k$ 次齐次的，则 $\varphi(\bar F)$ 是 $P$ 上的 $2k$ 次微分形式。

一个应当记住的典型例子是某个表示 $\rho$ 下的迹：
\begin{equation}
  \varphi_{k,\rho}(\bar F) := \operatorname{Tr}_\rho \rho(\bar F)^k
  \tag{I.10}
  \label{eq:I.10}
\end{equation}
（下文我们将经常略去此记号中的 $\rho$。）

\textbf{命题}：形式 $\varphi(\bar F)=\pi^*(\bar\varphi)$ 是基本形式。

\textbf{证明} 显然 $\varphi(\bar F)$ 是水平的，即 $\iota(\xi(X))\varphi(\bar F)=0$，因为 $\iota(\xi(X))(\bar F)=0$。此外，
\begin{equation}
  R_g^*(\varphi(\bar F)) = \varphi(R_g^*(\bar F)) = \varphi(g\bar Fg^{-1}) = \varphi(\bar F)
  \tag{I.11}
  \label{eq:I.11}
\end{equation}
因而它沿纤维方向不变。\ $\clubsuit$

\textbf{命题}：$\bar\varphi(\bar F)$ 是底空间 $M$ 上的闭微分形式。而且它是自然的。

\textbf{证明} 注意：
\begin{equation}
  d\varphi(\bar F,\cdots,\bar F) = k\,\varphi(d\bar F,\bar F,\cdots,\bar F)
  \tag{I.12}
  \label{eq:I.12}
\end{equation}
但现在我们使用 Bianchi 恒等式
\begin{equation}
  D\bar F = d\bar F + [\Theta,\bar F] = 0
  \tag{I.13}
  \label{eq:I.13}
\end{equation}
现在注意 $\varphi([\Theta,\bar F],\bar F,\cdots,\bar F)=0$，因为 $\varphi$ 是 $Ad$-不变的。（回想 $\Theta$ 取值于 $\mathfrak{g}$，并利用 $Ad$-不变性的无穷小形式。）于是（在 DeRham 复形上！）$\pi^*$ 没有核，故 $d\bar\varphi(\bar F)=0$。\ $\clubsuit$

% [Tutor Note: (I.12) 的关键是对 $k$ 个相同的偶次因子 $\bar F$（每个 2 次）用 Leibniz 法则：$k$ 个因子各贡献一项、没有符号因子（因 $\bar F$ 是偶次形式），故系数为 $k$；再代入 (I.13) 得 $d\bar F=-[\Theta,\bar F]$。而 $\varphi([\Theta,\bar F],\bar F,\dots,\bar F)=0$ 用到的是 $Ad$-不变性的\emph{无穷小}形式：$\sum_{i}\varphi(\bar F,\dots,[\Theta,\bar F],\dots,\bar F)=0$，即 $\varphi$ 在 $\mathfrak{g}$ 上 $ad(\mathfrak{g})$-不变。此处作者只说「利用 $Ad$-不变性的无穷小形式」，跳过了这层等价性。]

为证明自然性：

注意若 $f: P_1\to P_2$ 是覆盖 $\bar f: M_1\to M_2$ 的丛映射，且 $\Theta_1=f^*(\Theta_2)$ 是拉回联络，则 $\varphi(\bar F_1)=f^*(\varphi(\bar F_2))$，从而 $\bar\varphi(\bar F_1)=\bar f^*(\bar\varphi(\bar F_2))$。特别地，$\bar\varphi(\bar F)$ 的上同调类只依赖于 $(P,\Theta)$ 的同构类。\ $\clubsuit$

最后注意，若 $\varphi$ 是 $k$ 次的，则 $\bar\varphi(\bar F)\in\Omega^{2k}(M)$ 是 $2k$ 次形式。此外，对称多项式的乘积对应上同调类的乘积。于是我们有如下定义：

\textbf{定义}：Chern-Weil 同态是分次环的同态
\begin{equation}
  I^*(G)\to H^*_{DR}(M)
  \tag{I.14}
  \label{eq:I.14}
\end{equation}
它把 $\varphi$ 送到 $[\bar\varphi(\bar F)]$。

\textbf{注}：

\begin{enumerate}
  \item 考虑 Hopf 纤维化上的自然联络。由于 $u(1)$ 是阿贝尔的，Lie 代数上任何多项式都是不变的。场强 $\bar F$ 是基本形式，它在 $S^2$ 底空间上有非平凡的积分。{\small\itshape ［原文标记：$\clubsuit$需要解释$\clubsuit$］}
  % [原文待办标记，原文 p.215 右侧边注；原文写作 "NEED TO EXPLAIN"]
  \item 理解「该形式是基本的」这一事实的另一条途径，是完全在底空间上用局部坐标卡来做。在每个坐标卡 $\mathcal{U}_\alpha$ 中，$\operatorname{Tr}F_\alpha^k$ 是规范不变的。因此，跨过坐标卡边界比较联络得 $F_\alpha=g_{\alpha\beta}F_\beta g_{\alpha\beta}^{-1}$，从而 $\operatorname{Tr}F_\alpha^k$ 实际上是 $M$ 上整体良定义的 $2k$ 形式。在下面的工作中，我们有时会直接用定义在底空间上的 $F$ 来工作。回想在局部坐标下 $\Theta=g^{-1}dg+g^{-1}Ag$，而 $\bar F=g^{-1}Fg$。
  \item 更一般地，任何规范不变的量，例如由联络构成的 Yang-Mills 作用量密度 $\operatorname{Tr}F*F$，都会在底空间上整体良定义。
\end{enumerate}

% [Tutor Note: 第 2 条中原文写作 "path boundaries"，从上下文及第 219 页练习中的 "patch boundaries" 判断应为 "patch boundaries"（坐标卡边界）；此处按后者译出。$F_\alpha=g_{\alpha\beta}F_\beta g_{\alpha\beta}^{-1}$ 之所以成立，是因为在重叠区 $\Theta_\alpha=g_{\alpha\beta}^{-1}\Theta_\beta g_{\alpha\beta}+g_{\alpha\beta}^{-1}dg_{\alpha\beta}$ 中规范势的那一项恰好抵消曲率，曲率只作张量式变换。]

\paragraph{练习}
设 $\alpha,\beta$ 是取值于 $\mathfrak{g}$ 的 $j$ 次与 $k$ 次微分形式。证明
\begin{equation}
  \boxed{\,d\operatorname{Tr}(\alpha\wedge\beta)
  = \operatorname{Tr}\left[D_A(\alpha)\wedge\beta+(-1)^k\alpha\wedge D_A(\beta)\right]\,}
  \tag{I.15}
  \label{eq:I.15}
\end{equation}
其中
\begin{equation}
\begin{aligned}
  D_A\alpha &= d\alpha+A\alpha-(-1)^k\alpha A\\
  D_A\beta &= d\beta+A\beta-(-1)^j\beta A
\end{aligned}
  \tag{I.16}
  \label{eq:I.16}
\end{equation}

利用此式与 Bianchi 恒等式，直接验证
\begin{equation}
  d\operatorname{Tr}F^k = k\operatorname{Tr}(D_AF)F^{k-1} = 0
  \tag{I.17}
  \label{eq:I.17}
\end{equation}

\subsection{对联络的依赖：特征类的构造} % 原文: I.5 Dependence on connection: Construction of characteristic classes

现在来看 Weil 同态的一个关键性质：它的像实际上并不依赖于 $P$ 上联络的选取。

下面我们证明上同调类 $[\bar\varphi(\bar F)]$ 实际上与 $P$ 上联络的选取无关。

让我们从主要的特殊情形开始，并在底空间上局部地进行。于是我们考察基本形式 $\operatorname{Tr}F^n$ 对联络选取的依赖。回想在无穷小变分下：
\begin{equation}
  \delta F = D_A(\delta A)
  \tag{I.18}
  \label{eq:I.18}
\end{equation}

因此
\begin{equation}
\begin{aligned}
  \delta\operatorname{Tr}F^n
  &= n\operatorname{Tr}(\delta F)F^{n-1}\\
  &= n\operatorname{Tr}\left(D_A(\delta A)F^{n-1}\right)\\
  &= n\,d\left(\operatorname{Tr}\delta A\,F^{n-1}\right)
\end{aligned}
  \tag{I.19}
  \label{eq:I.19}
\end{equation}
其中我们用了 Bianchi 恒等式 $D_AF=0$。

更形式化地，若我们有一条\emph{线性}的联络路径：$A(t)=(1-t)A_0+tA_1=A_0+t\alpha$，$\alpha\in\Omega^1(M;\mathrm{ad}P)$，则
\begin{equation}
\begin{aligned}
  \operatorname{Tr}F_1^n-\operatorname{Tr}F_0^n
  &= \int_0^1 dt\,\frac{d}{dt}\operatorname{Tr}F(A(t))^n\\
  &= n\int_0^1 dt\,\operatorname{Tr}D_{A(t)}(\alpha)F(A(t))^{n-1}\\
  &= d\left[n\int_0^1 dt\,\operatorname{Tr}\alpha\,F(A(t))^{n-1}\right]\\
  &= d\left[n\int_0^1 dt\,\operatorname{Tr}\alpha\left(F_0+tD_{A_0}\alpha+t^2\alpha^2\right)^{n-1}\right]
\end{aligned}
  \tag{I.20}
  \label{eq:I.20}
\end{equation}

注意表达式
\begin{equation}
  n\int_0^1 dt\,\operatorname{Tr}\alpha\left(F_0+tD_{A_0}\alpha+t^2\alpha^2\right)^{n-1}
  \tag{I.21}
  \label{eq:I.21}
\end{equation}
是\emph{整体良定义}的，因为它只涉及张量性的量。因此，$M$ 上闭且整体良定义的形式 $\operatorname{Tr}F(A)^n$ 的 DeRham 上同调类与联络无关。

% [Tutor Note: (I.20) 第三行到第四行的替换是极易出错的一步：需先写出 $F(A(t))=d(A_0+t\alpha)+(A_0+t\alpha)^2=F_0+tD_{A_0}\alpha+t^2\alpha^2$，其中 $D_{A_0}\alpha=d\alpha+[A_0,\alpha]$；再把 $D_{A(t)}(\alpha)=D_{A_0}\alpha+2t\alpha^2$ 代回上一行，用 $\operatorname{Tr}$ 的循环性与 $D_{A_0}$ 下分部积分（$\int_0^1 dt\,d/dt(\cdot)$ 与 $\operatorname{Tr}\alpha\,d(\cdot)$ 的对易）把 $d$ 提到积分号外。作者只写了结果；核对时可用 $n=1$（此时 $F_0+tD_{A_0}\alpha+t^2\alpha^2$ 的零次幂为 $1$，退化为 $\operatorname{Tr}\alpha$）与 $n=2$ 两个情形验证系数。]

同样的运算对任意 $n$ 次不变多项式都成立：
\begin{equation}
\begin{aligned}
  \frac{d}{dt}\varphi(F(A(t)))
  &= n\,\varphi(D_{A(t)}\alpha,F(A(t)),\dots,F(A(t)))\\
  &= n\,d\Big(\varphi(\alpha,F(A(t)),\dots,F(A(t)))\Big)
\end{aligned}
  \tag{I.22}
  \label{eq:I.22}
\end{equation}
积分给出
\begin{equation}
  \varphi(F(A_1))-\varphi(F(A_0))
  = d\left[n\int_0^1\varphi(\alpha,F(A(t)),\dots,F(A(t)))\right]
  \tag{I.23}
  \label{eq:I.23}
\end{equation}
从而上同调类不变。

正如我们已经证明的，$[\bar\varphi(F)]$ 是自然的；因此，给定不变多项式 $\varphi$，我们可以按如下规则产生一个特征类：
\begin{equation}
  \xi: P\to[\bar\varphi(F)]
  \tag{I.24}
  \label{eq:I.24}
\end{equation}

其中我们可以使用 $P$ 上的任何联络。也就是说，Chern-Weil 同态与联络无关。

\textbf{注}：由此可知，闭链上的积分（包括 $M$ 闭时整个流形上的积分）都是「拓扑不变量」。这正是 Yang-Mills 作用量中许多拓扑项的来源。

\subsection{Borel 定理} % 原文: I.6 The Borel Theorem

现在我们可以把这一圈想法闭合起来：借助 Borel 定理，把不变多项式与分类空间的上同调联系起来：

\textbf{定理.} 设 $G$ 是紧致连通 Lie 群。则
\begin{equation}
  H^*(BG;\Lambda) = \left(H^*(BT,\Lambda)\right)^W
  \tag{I.25}
  \label{eq:I.25}
\end{equation}
其中 $W$ 是 Weyl 群，$T$ 是极大环面，$\Lambda$ 是任一使 $|W|$ 在其中可逆的环。

上标 $W$ 表示我们取 Weyl 群作用下的不变量。

若 $\Lambda=\mathbb{R}$，则我们可以用 Lie 代数上的 $Ad$-不变多项式来构造这个环：
\begin{equation}
  I^*(G) = S^*(\mathfrak{g}^*)^W
  \tag{I.26}
  \label{eq:I.26}
\end{equation}

于是

\textbf{定理}：若 $G$ 是紧致 Lie 群，则
\begin{equation}
  I^*(G)\cong H^*(BG;\mathbb{R})
  \tag{I.27}
  \label{eq:I.27}
\end{equation}

Borel 定理的参考文献：

\begin{enumerate}
  \item A. Borel, ``Topology of Lie Groups and Characteristic Classes,'' Bull. Amer. Math. Soc. \textbf{61}(1955) 397
  \item J. Dupont, \emph{Curvature and Characteristic Classes}, Lecture Notes in Mathematics \textbf{640} Springer 1978
\end{enumerate}

{\small\itshape ［原文标记：*******************\newline
讨论问题：若所有特征类都平凡，丛是否平凡？\newline
特征类的正确整归一化！！！！！\newline
*******************］}
% [原文待办标记，原文 p.217（印刷页），作者以星号行框住两句待办；原文为 "Discuss the question: If all the characteristic classes are trivial is the bundle trivial? / PROPER INTEGRAL NORMALIZATIONS OF THE CLASSES!!!!!"]

\subsection{次级特征类：Chern-Simons 形式} % 原文: I.7 Secondary characteristic classes: The Chern-Simons forms

让我们回到表达式：
\begin{equation}
  \varphi(F(A_1))-\varphi(F(A_0))
  = d\left[n\int_0^1 dt\,\varphi(\alpha,F(A(t)),\dots,F(A(t)))\right]
  \tag{I.28}
  \label{eq:I.28}
\end{equation}
它给出了 Chern-Weil 形式对联络选取的依赖。下面我们将更细致地考察那些使这一差别变为平凡的形式。为简单起见，我们把 $\varphi$ 取为 $\mathfrak{g}$ 的某个表示下的迹。

\emph{次级特征类}（secondary characteristic class）由使这一差别变为平凡的形式来定义：

\textbf{定义.} 两个联络之间的\emph{相对 Chern-Simons 形式}是如下形式：
\begin{equation}
  CS_{2n-1}(A_1,A_0) := n\int_0^1 dt\,\operatorname{Tr}\alpha
  \left(F_0+tD_{A_0}\alpha+t^2\alpha^2\right)^{n-1}
  \tag{I.29}
  \label{eq:I.29}
\end{equation}
其中 $A_1=A_0+\alpha$。

注意
\begin{equation}
  dCS_{2n-1}(A_1,A_0) = \operatorname{Tr}F_1^n-\operatorname{Tr}F_0^n
  \tag{I.30}
  \label{eq:I.30}
\end{equation}

\textbf{例}

\begin{equation}
  CS_1(A_1,A_0) = \operatorname{Tr}\alpha
  \tag{I.31}
  \label{eq:I.31}
\end{equation}

\begin{equation}
  CS_3(A_1,A_0) = \operatorname{Tr}\left(2\alpha F_0+\alpha D_{A_0}\alpha+\frac{2}{3}\alpha^3\right)
  \tag{I.32}
  \label{eq:I.32}
\end{equation}

\begin{equation}
\begin{aligned}
  CS_5(A_1,A_0) = \operatorname{Tr}\Big(&
    3\alpha F_0^2+\frac{3}{2}\alpha\{F_0,D_{A_0}\alpha\}+2F_0\alpha^3\\
  &+\alpha(D_{A_0}\alpha)^2+\frac{3}{2}\alpha^3D_{A_0}\alpha+\frac{3}{5}\alpha^5\Big)
\end{aligned}
  \tag{I.33}
  \label{eq:I.33}
\end{equation}

% [Tutor Note: 归一化最容易出错。(I.32) 中 $\alpha$ 的一次项系数 2、三次项系数 $2/3$ 都来自 (I.29) 对 $t$ 的积分：$n=2$ 时 $(F_0+tD_{A_0}\alpha+t^2\alpha^2)$ 的一次幂对 $t$ 积分给出 $F_0+\tfrac12 D_{A_0}\alpha+\tfrac13\alpha^2$，再乘以外面的 $n=2$ 即得 (I.32)。最省事的核对办法是取 $A_0=0$、$\alpha=A$，验证 $dCS_3=\operatorname{Tr}F^2$，即 $d\operatorname{Tr}(A\,dA+\tfrac23A^3)=\operatorname{Tr}(dA+A^2)^2$；若把系数写成 $1/3$ 或 $1/2$ 此式即不成立。]

\textbf{注}

\begin{enumerate}
  \item 通常 Chern-Simons 形式是用单个联络 $A$ 来表示的。当 $A$ 是平凡丛上的联络时，这样做是有意义的：此时可在上面所有公式中取 $A_0=0$、$A_1=A$。许多物理文献在这一点上非常马虎，当 $A$ 是非平凡联络时也写成 $CS(A)$。这种马虎会导致错误。务必小心！{\small\itshape ［原文标记：$\clubsuit$好吧。如果你取指数映射，那么对单个联络你\emph{确实}可以定义，而且这对后面各讲余下的内容很重要！！！$\clubsuit$］}
  % [原文待办标记，原文 p.218 右侧边注；原文写作 "Well. If you exponentiate, then you CAN define for a single connection, and this is important to the rest of the lectures !!!"]
  \item 话虽如此，确实有一种办法把 Chern-Simons 形式 $CS(A)$ 对单个联络定义为全空间上的良定义形式，即主丛 $P\to M$ 的全空间上的形式，因为联络形式在 $P$ 的全空间上是整体良定义的。更详细的讨论见 [22] 的附录 A。{\small\itshape ［原文标记：$\clubsuit$为什么 Dan 要用这个更复杂的论证？$\clubsuit$］}
  % [原文待办标记，原文 p.218 右侧边注；原文写作 "Why does Dan use this more complicated argument?"，Dan 指 D. Freed]
  \item 然后我们可以这样来定义 Chern-Simons 作用量：远离一个零测集选取截面 $s: M\to P$，并定义 $\int_M s^*CS(A)$。这将是不确定的，因为我们可以选取不同的截面，因此 $\int_M s^*CS(A)$ 只定义到模 $\mathbb{Z}$ 为止。
  \item 警告：在许多应用中，Chern-Simons 形式只被看作定义到恰当形式为止。注意 $CS\to CS+d(*)$ 不改变关键性质 \eqref{eq:I.20}。
\end{enumerate}

\paragraph{练习}
注意当 $M$ 不闭时，$\int_M\operatorname{Tr}F^n$ 依赖于联络。

\paragraph{练习}
证明：
\begin{equation}
  CS_3(A+\alpha+\beta,A)
  = CS_3(A+\alpha,A)+CS_3(A+\alpha+\beta,A+\alpha)-d(\operatorname{Tr}\alpha\beta)
  \tag{I.34}
  \label{eq:I.34}
\end{equation}
特别地：
\begin{equation}
  \frac{d}{dt}CS(A+\alpha(t),A)
  = 2\operatorname{Tr}\dot{\alpha}\big(F(A+\alpha)\big)-d(\operatorname{Tr}\alpha\dot{\alpha})
  \tag{I.35}
  \label{eq:I.35}
\end{equation}
这常被写成：
\begin{equation}
  \delta CS(A) = 2\operatorname{Tr}\delta A\wedge F(A)-d\big(\operatorname{Tr}(A\,\delta A)\big)
  \tag{I.36}
  \label{eq:I.36}
\end{equation}
这里假定丛是平凡的。

特别地，考虑 3 维流形 $M$ 上的主 $G$ 丛。「Chern-Simons 作用量」
\begin{equation}
  \int_M \operatorname{Tr}A\,dA+\frac{2}{3}A^3
  \tag{I.37}
  \label{eq:I.37}
\end{equation}
的变分是什么？

\paragraph{练习}
计算跨越坐标卡边界时 $CS_3(A_\alpha)-CS_3(A_\beta)$ 的变分。

\paragraph{练习}
证明 $CS(A_1^g,A_0^g)=CS(A_1,A_0)$。

现在让我们考察 Chern-Simons 形式本身对 $A$ 的依赖。

下面这个一般性的注记很有用：假设我们有一族联络 $A(s,x)$，记该族所在的空间为 $\mathcal{S}$。在 $\Omega^{p,q}(M\times\mathcal{S})$ 上引入双分次是有用的，使得 $(p,q)$ 形式属于 $\Omega^p(M)\otimes\Omega^q(\mathcal{S})$。我们可以在 $\pi^*P\to M\times\mathcal{S}$ 上构造一个联络 $\mathcal{A}(s,x)$，使其协变导数为
\begin{equation}
  d_M+d_{\mathcal{S}}+A(s,x)
  \tag{I.38}
  \label{eq:I.38}
\end{equation}
注意该联络形式只有 $(1,0)$ 型分量。这是所谓的「万有联络」（universal connection）的一个例子。（关于这一点更多内容，见下文。）设 $\mathcal{F}$ 是这个万有联络的曲率。它具有 $(2,0)$ 型与 $(1,1)$ 型分量。具体地
\begin{equation}
  \mathcal{F} = F(\mathcal{A}) = F(A)+d_{\mathcal{S}}A(s,x)
  \tag{I.39}
  \label{eq:I.39}
\end{equation}
现在 $\operatorname{Tr}\mathcal{F}^n$ 在全空间上是闭的，
\begin{equation}
  (d_M+d_{\mathcal{S}})\operatorname{Tr}\mathcal{F}^n = 0,
  \tag{I.40}
  \label{eq:I.40}
\end{equation}
于是若按 $(p,q)$ 型分解：
\begin{equation}
  \operatorname{Tr}\mathcal{F}^n = \omega^{n,0}+\omega^{n-1,1}+\dots+\omega^{0,n}
  \tag{I.41}
  \label{eq:I.41}
\end{equation}
并令 $d=d_M$、$\delta=d_{\mathcal{S}}$，我们就得到「下降方程」（descent equations）：
\begin{equation}
\begin{aligned}
  d\omega^{n,0} &= 0\\
  d\omega^{n-1,1}+\delta\omega^{n,0} &= 0\\
  d\omega^{n-2,2}+\delta\omega^{n-1,1} &= 0\\
  &\ \ \vdots
\end{aligned}
  \tag{I.42}
  \label{eq:I.42}
\end{equation}

% [Tutor Note: 原 PDF 第 220 页的 Figure 37 为空白占位框，图形内容在原稿中缺失，仅存图注；此处按原图注翻译并保留图号，以免后文图号错位。]
\begin{figure}[htbp]
  \centering
  \MissingFigureNote
  \caption{$\mathcal{A}(P)$ 中由三个联络构成的三角形。（原文 Fig.~37，第 220 页；原 PDF 此图缺失）}
  \label{fig:37}
\end{figure}

现在让我们考虑联络的一个单形 $\triangle$：
\begin{equation}
  A_0+t_1\alpha_1+t_2\alpha_2
  \tag{I.43}
  \label{eq:I.43}
\end{equation}
其中 $0\le t_2\le t_1\le 1$，如图 \ref{fig:37} 所示。

令 $A_1=A_0+\alpha_1$，$A_2=A_1+\alpha_2$。我们在 $\pi^*(P)\to M\times\triangle$ 上构造万有联络并计算曲率：
\begin{equation}
  F(\mathcal{A}) = F(A_0)+t_1D_{A_0}\alpha_1+t_2D_{A_0}\alpha_2
  +(t_1\alpha_1+t_2\alpha_2)^2+dt_1\alpha_1+dt_2\alpha_2
  \tag{I.44}
  \label{eq:I.44}
\end{equation}

利用下降方程可以写出
\begin{equation}
  \int_\triangle\delta\omega^{n-1,1}
  = -d\int_\triangle\omega^{n-2,2}
  = -d\int_\triangle\operatorname{Tr}F(\mathcal{A})^n
  \tag{I.45}
  \label{eq:I.45}
\end{equation}
另一方面，我们可以对三角形应用 Stokes 定理，写出
\begin{equation}
\begin{aligned}
  \int_\triangle\delta\omega^{n-1,1}
  &= \int_{\partial\triangle}\omega^{n-1,1}
  = \int_{\partial\triangle}\operatorname{Tr}F(\mathcal{A})^n\\
  &= CS_{2n-1}(A_1,A_0)+CS_{2n-1}(A_2,A_1)+CS_{2n-1}(A_0,A_2)
\end{aligned}
  \tag{I.46}
  \label{eq:I.46}
\end{equation}
把它们放在一起，我们就得到结果：
\begin{equation}
\begin{aligned}
  &CS_{2n-1}(A_0,A_1)+CS_{2n-1}(A_1,A_2)+CS_{2n-1}(A_2,A_0)\\
  &\quad = d_M\Bigg(\int_\triangle \operatorname{SymTr}\Big(dt_1\alpha_1,\,dt_2\alpha_2,\,
    \big(F(A_0)+t_1D_{A_0}\alpha_1+t_2D_{A_0}\alpha_2+(t_1\alpha_1+t_2\alpha_2)^2\big)^{n-2}\Big)\Bigg)
\end{aligned}
  \tag{I.47}
  \label{eq:I.47}
\end{equation}

于是，例如：
\begin{equation}
  CS_3(A_1,A_0)+CS_3(A_2,A_1)+CS_3(A_0,A_2) = d(\operatorname{Tr}\alpha_1\alpha_2)
  \tag{I.48}
  \label{eq:I.48}
\end{equation}
其中 $A_1=A_0+\alpha_1$，$A_2=A_1+\alpha_2$。

类似地：
\begin{equation}
\begin{aligned}
  &CS_5(A_1,A_0)+CS_5(A_2,A_1)+CS_5(A_0,A_2) =\\
  &\quad 3d\left[\int_\triangle \operatorname{Tr}
    \big(F+t_1D\alpha_1+t_2D\alpha_2+(t_1\alpha_1+t_2\alpha_2)^2\big)(\alpha_1\alpha_2-\alpha_2\alpha_1)\right]
\end{aligned}
  \tag{I.49}
  \label{eq:I.49}
\end{equation}

% [Tutor Note: 原 PDF 第 221 页的 Figure 38 为空白占位框，图形内容在原稿中缺失，仅存图注；此处按原图注翻译并保留图号，以免后文图号错位。]
\begin{figure}[htbp]
  \centering
  \MissingFigureNote
  \caption{$\mathcal{A}(P)$ 中用来计算 Chern-Simons 形式 $CS(A,0)$ 的规范依赖性的联络三角形。（原文 Fig.~38，第 221 页；原 PDF 此图缺失）}
  \label{fig:38}
\end{figure}

作为一个应用，设 $A$ 是可平凡化丛上的一个联络，于是存在基点联络 $A_0=0$。这时定义 $CS(A):=CS(A,0)$ 是有意义的。虽然 $dCS(A)$ 是规范不变的，$CS(A^g)$ 却不是规范不变的。

在上面的公式中，我们可以取 $A_1=g^{-1}dg$、$A_2=A^g$，其中 $A$ 是某个整体定义的联络 1-形式。注意这意味着我们有一条从 $A_1$ 到 $A_2$ 的线性路径 $A(t)=g^{-1}dg+t\,g^{-1}Ag$。由于一般地 $CS(A_1^g,A_0^g)=CS(A_1,A_0)$ 对任意一对联络成立，我们有
\begin{equation}
  CS_{2n-1}(A^g,g^{-1}dg) = CS_{2n-1}(A,0)
  \tag{I.50}
  \label{eq:I.50}
\end{equation}

因此，我们的恒等式变成：
\begin{equation}
  CS_{2n-1}(A^g)-CS_{2n-1}(A) = CS_{2n-1}(g^{-1}dg)-d_M(\alpha_{2n})
  \tag{I.51}
  \label{eq:I.51}
\end{equation}
其中
\begin{equation}
\begin{aligned}
  \alpha_{2n} = \int_\triangle \operatorname{SymTr}\Big[&
    dt_1\alpha_1,\,dt_2\alpha_2,\,
    \Big((t_1^2-t_1)\alpha_1^2+(t_2^2-t_2)\alpha_2^2\\
  &+(t_1t_2-t_2)[\alpha_1,\alpha_2]_+ + t_2g^{-1}F(A)g\Big)^{n-2}\Big]
\end{aligned}
  \tag{I.52}
  \label{eq:I.52}
\end{equation}
这里 $\alpha_1=g^{-1}dg$、$\alpha_2=g^{-1}Ag$，以及
\begin{equation}
  CS_{2n-1}(g^{-1}dg) = (-1)^{n-1}\frac{n!\,(n-1)!}{(2n-1)!}
  \operatorname{Tr}(g^{-1}dg)^{2n-1}
  \tag{I.53}
  \label{eq:I.53}
\end{equation}

% [Tutor Note: (I.53) 的系数可直接由 (I.29) 核对：取 $A_0=0$、$\alpha=g^{-1}dg$，此时 $\alpha$ 平坦（$d\alpha+\alpha^2=0$），故 $F(A(t))=t\,d\alpha+t^2\alpha^2=t(t-1)\alpha^2$，(I.29) 化为 $n\int_0^1 t^{n-1}(t-1)^{n-1}dt\,\operatorname{Tr}\alpha^{2n-1}=(-1)^{n-1}\frac{n!\,(n-1)!}{(2n-1)!}\operatorname{Tr}\alpha^{2n-1}$，正是 (I.53)。最容易抄错的是 Beta 积分 $B(n,n)=(n-1)!^2/(2n-1)!$ 给出的 $(n-1)!^2$ 与 $(-1)^{n-1}$；用 $n=2$ 检验：$2\int_0^1 t(t-1)dt=2\cdot(-\tfrac16)=-\tfrac13$，与 $(-1)^{1}\frac{2!\,1!}{3!}=-\tfrac13$ 一致。]

这个公式在讨论带 Chern-Simons 项的物理作用量、以及诸如反常的下降形式体系（descent formalism）与「反常流入」（anomaly inflow）等话题时很有用。见下文。

也可以用「Cartan 同伦算子」的形式体系来导出它。见

\begin{enumerate}
  \item Nakahara, sec. 11.5.3
  \item B. Zumino, Les Houches lectures
  \item Alvarez-Gaum\'e and Ginsparg, Ann. Phys.
\end{enumerate}

\paragraph{练习}
证明：若 $A$ 是丛 $P$（不必可平凡化）上的平坦联络，则
\begin{equation}
  CS(A^g,A) = (-1)^{n-1}\frac{n!\,(n-1)!}{(2n-1)!}\operatorname{Tr}(g^{-1}D_Ag)^{2n-1}
  \tag{I.54}
  \label{eq:I.54}
\end{equation}
其中 $g^{-1}D_Ag=g^{-1}(dg+Ag-gA)=A^g-A$。

\paragraph{练习（推广 Polyakov-Wiegmann 公式）}
在 \eqref{eq:I.51} 中取 $A=0^h$，导出 Polyakov-Wiegmann 公式及其各种推广。

\paragraph{练习}
设 $E=M\times\mathbb{C}^N$ 是平凡向量丛，$\mathcal{P}$ 与 $\mathcal{Q}$ 是到可能非平凡子丛的正交投影算子。

计算 $CS(\mathcal{P}d\oplus\mathcal{Q}d,d)$，其中 $d$ 是 $E$ 上的平凡联络。
